\documentclass[unicode,12pt,report]{ltjsarticle} 
\usepackage{amsmath, amssymb, amsthm}
\usepackage{mathrsfs}
\usepackage{tikz-cd}
\usetikzlibrary{calc}
\tikzset{curve/.style={settings={#1},to path={(\tikztostart)
    .. controls ($(\tikztostart)!\pv{pos}!(\tikztotarget)!\pv{height}!270:(\tikztotarget)$)
    and ($(\tikztostart)!1-\pv{pos}!(\tikztotarget)!\pv{height}!270:(\tikztotarget)$)
    .. (\tikztotarget)\tikztonodes}},
    settings/.code={\tikzset{quiver/.cd,#1}
        \def\pv##1{\pgfkeysvalueof{/tikz/quiver/##1}}},
    quiver/.cd,pos/.initial=0.35,height/.initial=0}
\usetikzlibrary{spath3}
\tikzset{between/.style n args={2}{/tikz/execute at end to={ \tikzset{spath/split at keep middle={current}{#1}{#2}} }}}
\usepackage[unicode, hidelinks]{hyperref} 
\usepackage{bbm}
\theoremstyle{definition} 
\newtheorem{definition}{定義}[section] 
\newtheorem{example}[definition]{例} 

\title{集合論}
\author{Soma Saito}
\date{\today}

\begin{document}
\maketitle
\tableofcontents
\clearpage

\section{はじめに}
このPDFでは最低限の集合論の基礎について解説する。
主に集合や写像の定義、用いられる記号や基本的な数学的操作に慣れ、数学に関する文章を一人でも読めるようになれることを目的としている。

\section{集合と写像}
\subsection{集合}
ほとんどの人がしっているとは思うが集合とは、ある対象がその集まりに入っているかどうかを客観的に、明確に判断できるものを集合(Set)という。
具体的には、自然数全体や、正$n$角形全体の集まりはもちろん集合である。しかし、「大きな円の集まり」や、「小さな有理数の集まり」は集合とはならない。
集合を構成する一つ一つの要素、つまり集合の中身のことを元と呼ぶ。また、ある集合$A$に元$x$が含まれていることを$x \in A$と書き、そうでないことを$x \notin  A$と記述する。
\subsection{集合の表し方}
集合には二つの表し方がある。外延的記法と内包的記法である。
外延的記法とは、その集合の元をすべて具体的に書き並べる方法である。例えば$A = \{1, 2, 3\}$など。
内包的記法とは、その元$x$が満たすべき条件を$P(x)$として、$A = \{x | P(x)\} (条件Pを満たすxの集まり)$のように書く。
例えば、$A = \{x | xは偶数\} = \{2, 4, 6, ...\}$である。
また、よく使う集合には特別な記号が割り当てられることもある。これらは宣言してから用いることで、説明を省くときに使う。たとえば、このPDF上では、自然数全体を$\mathbb{N}$とする。同様に整数や有理数も、
$\mathbb{Z}, \mathbb{Q}, \mathbb{R}$,さらに複素数の場合は$\mathbb{C}$である。
\subsection{部分集合}
\begin{definition}[部分集合]
    任意の$x \in A$に対して、$x \in B$が成り立つとき、$A$は$B$の部分集合であり、
    $A\subseteq B$と書く。
\end{definition}

\begin{definition}[集合の相等]
    $A, B$を集合とする。$A$と$B$が等しい、すなわち$A = B$であるというのは、$A \subseteq B$かつ$B \subseteq A$が成り立つことである。
\end{definition}

以下ではいくつかの集合の構成方法を紹介する。
\subsection{直積集合}
ここでは直積集合$A \times B$について紹介する。
直積集合とは以下のようなものである。
$$A \times B = \{(a, b) \mid a \in A かつ b \in B\}$$
つまり二つ組$a, b$のことであり、それぞれの元を一つずつ取り出している場合である。
さらに、$A \times A$のことを$A^2$と書くこともある。
実数全体に関して直積をとることで、いわゆるユークリッド空間(二次元平面、三次元空間)を構成できる。

似ている記号として、$A$から$B$への関数全体、関数空間として$B^A$などがあるがここでは扱わない。

\subsection{冪集合}
集合$X$に対して、冪集合は、$\mathscr{P}(X)$と表され、以下で定義される。
$$\mathscr{P}(X) = \{A \mid Xの部分集合\}$$
つまり、$X$のすべての部分集合が入っているような集合である。
例えば、$X = {a, b, c}$のとき、$\mathscr{P}(X) = \{\{a, b, c\}, \{a, b\}, \{b, c\}, \{a, c\}, \{a\}, \{b\}, \{c\}\}$
また、冪集合のことを$2^X$と書くこともある。これは、元の個数に由来している。

\subsection{写像}
ここでは写像について述べる。
一言で簡単に言えば、集合$A$から集合$B$への写像$f$とは、$A$の各要素を$B$の何らかの要素に対応づける規則のことである。数学的には、集合$A$のすべての元に対して、
集合$B$の要素がただ一つ定まる規則$f$を$A$から$B$への写像といい、次のように書く。$$f:A \to B$$ \\
$$x \mapsto f(x)$$

ここで、写像$f$の定義域とは集合$A$のことであり、終域とは集合$B$のことである。値とは、$x \in A$に対して定まる$f(x) \in B$のことである。
次に、重要な概念である全射、単射、全単射について述べる。
\begin{definition}[全射、単射、全単射]
    写像$f:A \to B$が全射(Surjctive)であるとは、集合$B$のすべての要素$b \in B$に対して、$f(x) = b$となる$x \in A$が存在することである。
    単射(Injective)であるとは、$x_1, x_2 \in A$のとき$x_1 \neq x_2$ならば$f(x_1) \neq f(x_2)$となることである。対偶を用いることも多い。
    $f$が全射かつ単射のとき、これを全単射と呼ぶ。
\end{definition}
\begin{example}
    $f:\mathbb{Z} \to \mathbb{Z}, f(n) \mapsto 2n$とすると、これは単射である。
    $f:\mathbb{R} \to [0, \inf), f(x) = x^2$とすると全射となるが、単射ではない。
    $f:\mathbb{R} \to \mathbb{R}, f(x) = 2x + 1$とすると、これは全単射である。二次元平面におけるグラフを考えるとわかる。
\end{example}
\subsection{逆写像}
逆写像について述べる。
$f:A \to B, x \mapsto f(x)$が写像であるとき、\\
$f^{-1}: B \to A, f(x) \mapsto x$を逆写像という。
つまり、写像の行先から元の場所に戻るような対応関係を逆写像と呼ぶ。写像の定義から明らかではあるが、逆写像は一般には写像とはならない。
写像$f$が全単射であるとき、またその時に限り逆写像は写像となり、かつ全単射である。


\section{同値関係と商集合}
集合の元に対して、何らかの関係を用いてグループ分けをすることを考える。今回は、整数における通常の割り算のあまりのみについて取り扱うが、一般の集合と一般の関係についてこれは定義できることに注意する。
\begin{example}
    整数全体$\mathbb{Z}$を「$3$で割ったあまり」についてグループ分けする。
    \begin{itemize}
        \item あまり0のグループ:$\{..., -3, 0, 3, 6, ...\}$
        \item あまり1のグループ:$\{..., -2, 1, 4, 7, ...\}$
        \item あまり2のグループ:$\{..., -1, 2, 5, 8, ...\}$
    \end{itemize}
    ここで、$a, b \in \mathbb{Z}$が、$3$で割ったとき同じグループに属しているとき、すなわちあまりが等しいとき、$a \equiv b (mod 3)$とかくこととする。
\end{example}
このように、集合に対して$equiv$という「関係」を定義することで、それぞれの要素をグループ分けできる。


\subsection{同値関係}
次に一般の同値関係について述べる。
記号$\sim $が集合$X$上の同値関係であるとは、以下の三つを満たすことをいう。
\begin{itemize}
    \item [反射律] 任意の$x \in X$に対して、$x \sim x$
    \item [推移律] $x, y, z \in X$に対して、$x \sim y,  y \sim z$ならば$x \sim z$
    \item [対称律] $x, y \in X$に対して、$x \sim y$ならば$y \sim x$
\end{itemize}

このように同値関係を用いて各元を同じものとみなして分割することはよくある。
\subsection{同値類}
ここからは再びmodの例に戻って話を進める。
$a \in \mathbb{Z}$と同値関係にある要素をすべて集めた集合を、$a$の同値類と呼び、$[a]$とかく。
つまり、
\begin{itemize}
    \item $[0] = \{..., -3, 0, 3, 6, ...\}$
    \item $[1] = \{..., -2, 1, 4, 7, ...\}$
    \item $[2] = \{..., -1, 2, 5, 8, ...\}$
\end{itemize}
である。この場合の$0, 1, 2$を代表元と呼ぶ。ただし代表元の取り方には任意性があることに注意する。

\subsection{商集合}
\begin{definition}
    集合$A$を同値関係$\sim$で割った同値類の集まりを商集合といい、$A/\sim$とかく。
\end{definition}

\begin{example}
    mod 3の例では、$$\mathbb{Z}/3\mathbb{Z} = \{[0], [1], [2]\}$$
    となる。
\end{example}
同値関係を定義することで、集合の直和分割が自然に定まる。

\section{論理記号}
ここでは、数学に関する文章に出てくる記号をいくつか紹介する。
\begin{itemize}
    \item $\exists$(存在量化):「~が存在する」「ある~に対して」
    \item $\forall$(全称量化):「すべての」「任意の」
    \item $\wedge$ かつ
    \item $\vee $ または
    \item $\lnot $否定
    \item $\Rightarrow$ならば
    \item $\Leftrightarrow $ 同値、必要十分条件、iff \footnote{if and only ifの略。～のとき、かつその時に限り、という意味。}
\end{itemize}

\begin{example}[順番で意味が変わる話]
    $$\forall x \in \mathbb{R}, \exists y \in \mathbb{R} s.t. x + y = 0$$
    $$\exists y \in \mathbb{R}, \forall x \in \mathbb{R} s.t. x + y = 0$$   
    \footnote{s.t.はsuch thatの略。～となるような、という意味。}
    上の文は、すべての$x$に対して、ある$y$が存在して、足し算すると$0$になる。(正しい)\\
    下の文は、ある最強の$y$が存在して、すべての$x$に対してそれを足し算すると$0$になる。(そんなyはいない)
\end{example}

\section{おわりに}
ここでは簡単な集合論の基礎を解説した。追加すべき事項や数学的な不備があればぜひ連絡してください。
\end{document}